\documentclass[11pt,a4paper]{article}

\usepackage[utf8]{inputenc}
\usepackage[a4paper, margin=1in, headheight=14pt]{geometry}
\usepackage{lmodern}
\usepackage{microtype}
\usepackage{titlesec}
\usepackage{fancyhdr}
\usepackage{lastpage}
\usepackage{booktabs}
\usepackage{tabularx}
\usepackage{array}
\usepackage{graphicx}
\usepackage{float}
\usepackage[table]{xcolor}
\usepackage{enumitem}
\usepackage{listings}
\usepackage{amsmath}
\usepackage{amssymb}
\usepackage{pgfgantt}
\usepackage[hidelinks, pdfencoding=auto]{hyperref}

% Define Palette
\definecolor{primary}{RGB}{20, 45, 85}      % Deep Academic Navy
\definecolor{secondary}{RGB}{35, 110, 165}  % Slate Blue
\definecolor{accent}{RGB}{195, 68, 38}      % Warm Terracotta Accent
\definecolor{darkgray}{RGB}{50, 50, 50}
\definecolor{lightgray}{RGB}{245, 247, 250}
\definecolor{codebg}{RGB}{248, 249, 251}
\definecolor{bordercolor}{RGB}{215, 222, 232}
\definecolor{successgreen}{RGB}{34, 139, 34}

% Setup Titlesec
\titleformat{\section}{\Large\bfseries\color{primary}}{\thesection}{0.6em}{}[\titlerule]
\titleformat{\subsection}{\large\bfseries\color{secondary}}{\thesubsection}{0.6em}{}
\titleformat{\subsubsection}{\normalsize\bfseries\color{darkgray}}{\thesubsubsection}{0.6em}{}
\titlespacing*{\section}{0pt}{1.4em}{0.6em}
\titlespacing*{\subsection}{0pt}{1.1em}{0.4em}
\titlespacing*{\subsubsection}{0pt}{0.8em}{0.3em}

% Setup Headers & Footers
\pagestyle{fancy}
\fancyhf{}
\fancyhead[L]{\small\textsf{\color{darkgray}UCS503P: Software Engineering --- Prototype Stage Report}}
\fancyhead[R]{\small\textsf{\color{darkgray}CMMS Platform}}
\fancyfoot[L]{\small\textsf{\color{darkgray}Thapar Institute of Engineering \& Technology}}
\fancyfoot[R]{\small\textsf{\color{darkgray}Page \thepage\ of \pageref{LastPage}}}
\renewcommand{\headrulewidth}{0.4pt}
\renewcommand{\footrulewidth}{0.4pt}

% Code Listing Settings
\lstset{
  backgroundcolor=\color{codebg},
  basicstyle=\ttfamily\footnotesize\color{darkgray},
  breaklines=true,
  captionpos=b,
  frame=single,
  rulecolor=\color{bordercolor},
  keywordstyle=\color{secondary}\bfseries,
  stringstyle=\color{accent},
  commentstyle=\color{gray}\itshape,
  showstringspaces=false,
  tabsize=2
}

% Hyperref Setup
\hypersetup{
  colorlinks=true,
  linkcolor=primary,
  citecolor=secondary,
  urlcolor=secondary,
  pdftitle={College Mentorship Management System (CMMS) - Prototype Stage Report},
  pdfauthor={Aaditya Bansal, Jessica, Harshveer Singh}
}

\begin{document}

% ============================================================================
% TITLE PAGE
% ============================================================================
\begin{titlepage}
  \centering
  \vspace*{0.8cm}
  
  {\scshape\Large Thapar Institute of Engineering and Technology \par}
  \vspace{0.2cm}
  {\scshape\large Department of Computer Science and Engineering \par}
  \vspace{0.1cm}
  {\small Patiala, Punjab, India \par}
  
  \vspace{1.5cm}
  \rule{\linewidth}{1.5pt} \\[0.4cm]
  {\huge\bfseries\color{primary} College Mentorship Management System \\[0.2cm] (CMMS) \par}
  \vspace{0.3cm}
  {\large\bfseries\color{secondary} A Digital Full-Stack Web Platform for Faculty--Student Mentorship \par}
  \vspace{0.2cm}
  {\normalsize\itshape Prototype Stage Report \& Engineering Specification \par}
  \rule{\linewidth}{1.5pt}
  
  \vspace{1.5cm}
  
  \begin{minipage}[t]{0.48\textwidth}
    \textbf{\large Submitted by:}\\[0.4cm]
    \textbf{Aaditya Bansal} \\
    Roll No: 1024030768 \\
    \textit{Backend \& System Integration Lead} \\[0.3cm]
    \textbf{Jessica} \\
    Roll No: 1024030761 \\
    \textit{Frontend UI/UX Lead} \\[0.3cm]
    \textbf{Harshveer Singh} \\
    Roll No: 1024030776 \\
    \textit{Database \& Conflict Logic Lead}
  \end{minipage}
  \hfill
  \begin{minipage}[t]{0.48\textwidth}
    \raggedleft
    \textbf{\large Submitted to:}\\[0.4cm]
    \textbf{Jeelani Asif} \\
    Course Instructor / Evaluator \\[0.4cm]
    \textbf{Course Code:}\\[0.1cm]
    UCS503P: Software Engineering (Project) \\[0.4cm]
    \textbf{Academic Session:}\\[0.1cm]
    2026--2027 (Odd Semester)
  \end{minipage}
  
  \vfill
  \rule{0.6\linewidth}{0.4pt} \\
  {\small\textsf{September 2026}}
\end{titlepage}

\newpage
\pagenumbering{roman}
\tableofcontents
\vspace{0.8cm}
\listoffigures
\vspace{0.8cm}
\listoftables
\newpage
\pagenumbering{arabic}

% ============================================================================
% SECTION 1: EXECUTIVE SUMMARY & ABSTRACT
% ============================================================================
\section{Executive Summary \& Abstract}
Institutional mentorship in higher education is universally acknowledged as a critical catalyst for undergraduate retention, academic excellence, and student career readiness. Despite its paramount significance, faculty--student mentorship across collegiate departments remains largely manual, fragmented, and ad-hoc---relying on uncoordinated email chains, informal walk-ins during unmonitored office hours, and unrecorded interactions. This lack of structure leads to scheduling collisions, unequal faculty advising loads, zero longitudinal tracking of student difficulties, and zero administrative visibility.

The \textbf{College Mentorship Management System (CMMS)} addresses this institutional deficit by providing an end-to-end, multi-tier web application built upon the modern MERN architecture (React 18 + Vite, Node.js + Express 5, and MongoDB Atlas). This document represents the \textbf{Prototype Stage Report} for CMMS, marking the transition from theoretical system proposal to a fully functional, testable software implementation.

\begin{figure}[H]
\centering
\begin{tabular}{|p{0.95\textwidth}|}
\hline
\rowcolor{lightgray} \textbf{\color{primary}Summary of Prototype Accomplishments} \\
\hline
\begin{itemize}[leftmargin=1.4em, itemsep=0.15em]
  \item \textbf{Decoupled Three-Tier Architecture:} Complete operational separation of concerns between React SPA client, RESTful Express application server, and MongoDB document database.
  \item \textbf{Multi-Role Security (RBAC):} Secure token-based authentication (JSON Web Tokens and 10-round bcrypt salted password hashing) with role-guarded route trees for Students, Faculty Mentors, and Department Administrators.
  \item \textbf{Intra-Departmental Mentee Allocation Engine:} Backend algorithm matching students to faculty mentors within corresponding academic departments and focus areas, supported by administrative reassignment tools.
  \item \textbf{Conflict-Free Office Hour Booking:} Server-side slot publishing and collision detection ensuring zero double-booking across student schedules and faculty availability windows.
  \item \textbf{Post-Meeting Note Logging \& History:} Structured chronological audit trail of all faculty-authored meeting summaries and student ratings, enforcing accountability.
  \item \textbf{Three Role-Specific Dashboards:} Responsive, data-driven interfaces providing immediate contextual value to all three distinct user classes.
\end{itemize} \\
\hline
\end{tabular}
\caption{Prototype Stage Core Deliverables Summary}
\end{figure}

The prototype implementation comprises over 15 backend API endpoints and 12 modular frontend interface views. Verification against predefined, measurable acceptance criteria demonstrates 100\% allocation correctness, sub-200ms API response latency under simulated departmental loads, and complete data isolation between unrelated user entities.

% ============================================================================
% SECTION 2: PROBLEM STATEMENT, OBJECTIVES & PROTOTYPE SCOPE
% ============================================================================
\section{Problem Statement, Objectives \& Prototype Scope}

\subsection{Problem Context \& Current Gaps}
In contemporary collegiate engineering programs, students frequently encounter academic stagnation, curriculum disorientation, and career uncertainty. While formal institutional guidelines mandate that each student be allocated a faculty advisor, the actual day-to-day operation suffers from four systemic deficiencies:
\begin{enumerate}[leftmargin=1.4em, itemsep=0.3em]
  \item \textbf{Coordination Friction:} Meeting arrangements require protracted back-and-forth email negotiations. Misunderstandings regarding office availability lead to high cancellation rates and missed appointments.
  \item \textbf{Absence of Longitudinal Tracking:} When meetings do transpire, discussions are rarely documented in a centralized repository. When a student changes semesters or advisors, prior context is lost.
  \item \textbf{Advising Load Asymmetry:} In large departments, popular or senior professors inadvertently become overburdened with disproportionate numbers of advisees, while junior faculty remain underutilized.
  \item \textbf{Departmental Opacity:} Department heads and academic deans possess no real-time telemetry regarding mentorship health, unable to determine whether at-risk students are receiving timely guidance.
\end{enumerate}

\subsection{Formal Project Objectives}
The CMMS project is chartered to realize the following concrete engineering objectives:
\begin{itemize}[leftmargin=1.4em, itemsep=0.2em]
  \item \textbf{Automated \& Fair Allocation:} Implement a deterministic allocation engine that matches incoming students to department faculty based on stated focus areas and existing caseload limits.
  \item \textbf{Conflict-Free Scheduling System:} Provide an interactive time-slot booking mechanism that enforces concurrency checks, preventing overlapping bookings.
  \item \textbf{Auditable Interaction Logs:} Establish an immutable, timestamped record of every completed session, accompanied by faculty feedback and qualitative progress notes.
  \item \textbf{Departmental Telemetry Dashboard:} Deliver aggregated analytics (completion ratios, mentee distributions, active sessions) to department administrators without requiring manual inquiry.
\end{itemize}

\subsection{Scope Delimitation: Prototype Stage vs. Final Release}
To uphold agile development principles and maintain strict time-to-value heuristics, the software engineering lifecycle has been phased. Table~\ref{tab:scope} details the boundary of the current Prototype Stage relative to the Final Release.

\begin{table}[H]
\centering
\small
\begin{tabularx}{\textwidth}{lXcc}
\toprule
\textbf{Functional Module} & \textbf{Feature Description} & \textbf{Prototype Stage} & \textbf{Final Release} \\
\midrule
Authentication & Email/Password JWT Auth with Role Claims & \textbf{Implemented} & Implemented \\
Role-Based Access & Frontend route protection \& backend middleware & \textbf{Implemented} & Implemented \\
Mentor Allocation & Intra-department matching \& Admin Reassignment & \textbf{Implemented} & Implemented \\
Scheduling Engine & Availability slot creation \& conflict validation & \textbf{Implemented} & Implemented \\
Meeting Lifecycle & Scheduled $\to$ Completed $\to$ Cancelled transitions & \textbf{Implemented} & Implemented \\
Interaction Logs & Timestamped faculty meeting notes & \textbf{Implemented} & Implemented \\
Student Feedback & 1--5 star rating \& meeting comments & \textbf{Implemented} & Implemented \\
Admin Metrics & Mentee counts, active meetings, completion \% & \textbf{Implemented} & Implemented \\
Notifications & Synchronous in-app banners \& UI status alerts & \textbf{Implemented} & Implemented \\
Email / SMS Gateway & Automated SMTP email dispatch on booking & \textit{Deferred} & Implemented \\
Calendar Sync & Google Calendar / Microsoft Outlook iCal sync & \textit{Deferred} & Implemented \\
Dynamic Load Heuristic & Automated multi-variable optimization matching & \textit{Deferred} & Implemented \\
\bottomrule
\end{tabularx}
\caption{Prototype vs. Final System Scope Matrix}
\label{tab:scope}
\end{table}

% ============================================================================
% SECTION 3: SYSTEM ARCHITECTURE & TECH STACK
% ============================================================================
\section{Prototype System Architecture \& Tech Stack}

\subsection{Three-Tier Decoupled Architecture}
The system adheres to an industry-standard three-tier decoupled client-server architecture. This separation guarantees high modularity, testability, and horizontal scalability.

\begin{figure}[H]
\centering
\small
\begin{tabularx}{\textwidth}{|X|X|X|}
\hline
\rowcolor{lightgray} \centering\textbf{Tier 1: Client / Presentation} & \centering\textbf{Tier 2: Application / REST API} & \centering\textbf{Tier 3: Persistence Layer} \tabularnewline
\hline
\begin{itemize}[leftmargin=1.2em, itemsep=0.1em, topsep=0.2em]
  \item React 18 SPA with Vite
  \item React Router DOM v6 (RBAC)
  \item Context API (\texttt{AuthContext})
  \item Responsive CSS Design System
  \item Axios API Interceptor Client
\end{itemize} & 
\begin{itemize}[leftmargin=1.2em, itemsep=0.1em, topsep=0.2em]
  \item Node.js LTS \& Express.js 5
  \item JWT Bearer Authentication
  \item Bcryptjs Password Hashing
  \item Conflict-Free Booking Engine
  \item Async Error Bubble Pipeline
\end{itemize} & 
\begin{itemize}[leftmargin=1.2em, itemsep=0.1em, topsep=0.2em]
  \item MongoDB Atlas Cloud NoSQL
  \item Mongoose ODM 8 Schemas
  \item Strong Relational ObjectIds
  \item Compound Conflict Indexes
  \item Normalized \& Embedded Data
\end{itemize} \tabularnewline
\hline
\end{tabularx}
\caption{Three-Tier System Architecture Decomposition}
\label{fig:tiers}
\end{figure}

\subsection{Technology Stack Selection Rationale}
The technology stack choices were governed by the ``Engine Availability'' and ``Time-to-Value'' engineering heuristics:
\begin{itemize}[leftmargin=1.4em, itemsep=0.3em]
  \item \textbf{Frontend (React 18 + Vite):} React's component-driven paradigm enables granular UI state encapsulation. Vite provides sub-second hot-module replacement (HMR), drastically reducing frontend development and testing cycles.
  \item \textbf{Backend (Node.js + Express.js 5):} Node's asynchronous, event-driven I/O model handles concurrent HTTP connections with minimal memory overhead, perfectly suited for rapid request-response API gateways. Express 5 provides modernized Promise error handling natively.
  \item \textbf{Database (MongoDB Atlas with Mongoose):} The flexible document model naturally represents hierarchical academic records (e.g., student profiles with embedded meeting references and expertise arrays) without rigid multi-table SQL join overhead. Mongoose enforces schema validation and index constraints.
\end{itemize}

% ============================================================================
% SECTION 4: FORMAL SYSTEM DESIGN & ARCHITECTURAL MODELS
% ============================================================================
\section{Formal System Design \& Architectural Models}

\subsection{Use-Case Modeling}
The system boundary encompasses three primary human actors: \textbf{Student}, \textbf{Faculty Mentor}, and \textbf{Department Administrator}. Figure~\ref{fig:usecase} illustrates the actor interactions and system use-case relationships.

\begin{figure}[H]
\centering
\includegraphics[width=0.72\textwidth,height=0.25\textheight,keepaspectratio]{images/usecase.png}
\caption{CMMS System Use-Case Diagram (Actors and Use-Case Boundaries)}
\label{fig:usecase}
\end{figure}

\begin{itemize}[leftmargin=1.4em, itemsep=0.2em]
  \item \textbf{Student Actions:} Register/Authenticate, view assigned faculty mentor details, browse open availability slots, book conflict-free office hours, review historical meeting summaries, and submit post-session feedback.
  \item \textbf{Faculty Actions:} Authenticate, publish open office-hour time slots, view active advisee roster, manage session lifecycle (confirm or cancel appointments), and append confidential post-meeting notes.
  \item \textbf{Administrator Actions:} Access departmental macro-analytics, observe faculty mentee load distributions, filter engagement metrics, and execute manual mentee reassignments.
\end{itemize}

\subsection{Data Flow Modeling}

\subsubsection{Level 0 DFD (Context Diagram)}
The Context Diagram (Figure~\ref{fig:dfd0}) portrays the CMMS as a singular system process interacting with external entities.

\begin{figure}[H]
\centering
\includegraphics[width=0.72\textwidth,height=0.22\textheight,keepaspectratio]{images/dfd_level0.png}
\caption{Level 0 Data Flow Diagram: Context-Level Process Flow}
\label{fig:dfd0}
\end{figure}

\subsubsection{Level 1 DFD (Module Decomposition)}
The Level 1 DFD (Figure~\ref{fig:dfd1}) decomposes the single system into its four core operational subsystems: Mentee Allocation (1.0), Scheduling \& Booking (2.0), Meeting Logging (3.0), and Dashboards \& Analytics (4.0), mapping data flows to their underlying persistence stores.

\begin{figure}[H]
\centering
\includegraphics[width=0.75\textwidth,height=0.26\textheight,keepaspectratio]{images/dfd_level1.png}
\caption{Level 1 Data Flow Diagram: Core Subsystems and Data Stores}
\label{fig:dfd1}
\end{figure}

\subsubsection{Level 2 DFD (Process 2.0: Scheduling \& Booking Decomposition)}
Process 2.0 represents the highest complexity subsystem in the application. Its internal transaction flow is partitioned into four distinct sub-processes:
\begin{enumerate}[leftmargin=1.4em, itemsep=0.2em]
  \item \textbf{2.1 Create Availability Slots:} Faculty submits time windows; server persists them to the Scheduling Store (\texttt{D2}).
  \item \textbf{2.2 Fetch Available Slots:} Student queries open slots for their specific assigned mentor; system filters out booked intervals.
  \item \textbf{2.3 Validate Conflict-Free Booking:} The booking engine executes an atomic query ensuring the slot has not been concurrently reserved by another student.
  \item \textbf{2.4 Generate Booking Log:} The confirmed session is written to the Meeting Collection and dispatch notifications are generated.
\end{enumerate}

\subsection{Entity-Relationship Model \& Data Architecture}
Figure~\ref{fig:er} models the conceptual entity schema, showcasing primary foreign key relationships between the system's foundational collections.

\begin{figure}[H]
\centering
\includegraphics[width=0.80\textwidth,height=0.28\textheight,keepaspectratio]{images/er_diagram.png}
\caption{CMMS Entity-Relationship Diagram (ERD)}
\label{fig:er}
\end{figure}

% ============================================================================
% SECTION 5: DATABASE SCHEMA & ENTITY SPECIFICATIONS
% ============================================================================
\section{Database Schema \& Entity Specifications}
The prototype utilizes Mongoose ODM to enforce strict schema structures, data types, validation rules, and relational references across MongoDB documents.

\subsection{User Collection Schema (\texttt{User.js})}
The root entity for authentication and RBAC authorization:
\begin{table}[H]
\centering
\small
\begin{tabularx}{\textwidth}{llp{3cm}X}
\toprule
\textbf{Field Name} & \textbf{Data Type} & \textbf{Constraints} & \textbf{Description} \\
\midrule
\texttt{\_id} & ObjectId & Primary Key, Auto & Unique internal identifier \\
\texttt{name} & String & Required, Trimmed & Full name of the user \\
\texttt{email} & String & Required, Unique, Lower & University email credential \\
\texttt{password} & String & Required, MinLen 6 & 10-round bcrypt hash string \\
\texttt{role} & String & Enum: \texttt{student}, \texttt{faculty}, \texttt{admin} & System RBAC role descriptor \\
\texttt{createdAt} & Date & Auto Timestamp & Account creation timestamp \\
\bottomrule
\end{tabularx}
\caption{User Collection Schema Specification}
\end{table}

\subsection{Student Collection Schema (\texttt{Student.js})}
Captures academic standing and mentor affiliation:
\begin{table}[H]
\centering
\small
\begin{tabularx}{\textwidth}{llp{3cm}X}
\toprule
\textbf{Field Name} & \textbf{Data Type} & \textbf{Constraints} & \textbf{Description} \\
\midrule
\texttt{\_id} & ObjectId & Primary Key, Auto & Unique student record ID \\
\texttt{user} & ObjectId & Ref: \texttt{User}, Required & Pointer to parent auth credential \\
\texttt{rollNumber} & String & Required, Unique, Trim & Institutional registration number \\
\texttt{department} & String & Required, Trimmed & Academic department (e.g., CSE) \\
\texttt{semester} & Number & Integer, Min 1, Max 12 & Current active semester \\
\texttt{mentor} & ObjectId & Ref: \texttt{Mentor}, Optional & Currently assigned faculty advisor \\
\bottomrule
\end{tabularx}
\caption{Student Collection Schema Specification}
\end{table}

\subsection{Mentor Collection Schema (\texttt{Mentor.js})}
Stores professional specialization and office-hour windows:
\begin{table}[H]
\centering
\small
\begin{tabularx}{\textwidth}{llp{3cm}X}
\toprule
\textbf{Field Name} & \textbf{Data Type} & \textbf{Constraints} & \textbf{Description} \\
\midrule
\texttt{\_id} & ObjectId & Primary Key, Auto & Unique mentor record ID \\
\texttt{user} & ObjectId & Ref: \texttt{User}, Required & Pointer to parent auth credential \\
\texttt{department} & String & Required, Trimmed & Academic department affiliation \\
\texttt{expertise} & Array [String] & Optional & Technical or academic domains \\
\texttt{availability} & String & Formatted JSON / Text & Active office hour schedule description \\
\bottomrule
\end{tabularx}
\caption{Mentor Collection Schema Specification}
\end{table}

\subsection{Meeting Collection Schema (\texttt{Meeting.js})}
Manages the appointment lifecycle and interaction state machine:
\begin{table}[H]
\centering
\small
\begin{tabularx}{\textwidth}{llp{3cm}X}
\toprule
\textbf{Field Name} & \textbf{Data Type} & \textbf{Constraints} & \textbf{Description} \\
\midrule
\texttt{\_id} & ObjectId & Primary Key, Auto & Unique meeting transaction ID \\
\texttt{mentor} & ObjectId & Ref: \texttt{Mentor}, Required & Faculty advisor participating \\
\texttt{student} & ObjectId & Ref: \texttt{Student}, Required & Student advisee participating \\
\texttt{date} & Date & Required, Indexed & ISO Date-Time of the session \\
\texttt{agenda} & String & Required, Trimmed & Student-stated objective for session \\
\texttt{status} & String & Enum: \texttt{scheduled}, \texttt{completed}, \texttt{cancelled} & State of meeting lifecycle \\
\texttt{notes} & String & Optional, Trimmed & Immediate outcome summary note \\
\bottomrule
\end{tabularx}
\caption{Meeting Collection Schema Specification}
\end{table}

\subsection{Supplementary Collections (\texttt{Note}, \texttt{Assignment}, \texttt{Feedback})}
\begin{itemize}[leftmargin=1.4em, itemsep=0.2em]
  \item \textbf{Note Collection (\texttt{Note.js}):} Decoupled collection containing detailed historical records (\texttt{meetingId: ObjectId}, \texttt{content: String}, \texttt{createdAt: Date}), ensuring granular timestamped documentation.
  \item \textbf{Assignment Collection (\texttt{Assignment.js}):} Tracks longitudinal mentor-student allocations, recording \texttt{studentId}, \texttt{facultyId}, assignment status (\texttt{active}, \texttt{reassigned}, \texttt{completed}), and timestamps.
  \item \textbf{Feedback Collection (\texttt{Feedback.js}):} Captures post-meeting student evaluations (\texttt{meeting: ObjectId}, \texttt{student: ObjectId}, \texttt{rating: 1..5}, \texttt{comment: String}) for quality assurance metrics.
\end{itemize}

% ============================================================================
% SECTION 6: RESTFUL API SPECIFICATION
% ============================================================================
\section{RESTful API Specification}
The backend service exposes a clean, deterministic REST API structured around resource routes. All protected routes require a bearer token in the HTTP Authorization header: \texttt{Authorization: Bearer <JWT>}.

\subsection{API Route Catalog}
Table~\ref{tab:api} summarizes the primary endpoints implemented and tested in the prototype.

\begin{table}[H]
\centering
\scriptsize
\begin{tabularx}{\textwidth}{@{}p{1.0cm}>{\ttfamily\raggedright}p{4.6cm}p{1.3cm}>{\raggedright}p{2.8cm}X@{}}
\toprule
\textbf{Method} & \textbf{Route Path} & \textbf{Auth} & \textbf{Payload / Params} & \textbf{Primary Function} \\
\midrule
\texttt{GET} & /api/health & Public & None & Service liveness \& health probe \\
\texttt{POST} & /api/users & Public & \texttt{\{name, email, pwd, role\}} & User registration \& password hashing \\
\texttt{POST} & /api/users/login & Public & \texttt{\{email, password\}} & JWT generation \& credentials check \\
\texttt{GET} & /api/users & Admin & None & List all system user records \\
\midrule
\texttt{POST} & /api/mentors & Admin & \texttt{\{user, dept, expertise\}} & Create faculty mentor profile \\
\texttt{GET} & /api/mentors & Protected & None & Retrieve all active mentors \\
\texttt{GET} & /api/mentors/:id/mentees & Faculty & URL param \texttt{:id} & Retrieve students assigned to mentor \\
\texttt{PUT} & /api/mentors/:id/availability & Faculty & \texttt{\{availability\}} & Update office-hour available slots \\
\midrule
\texttt{POST} & /api/students & Protected & \texttt{\{user, rollNo, dept\}} & Create student profile \\
\texttt{GET} & /api/students & Protected & None & List all student profiles \\
\texttt{GET} & /api/students/:id/mentor & Student & URL param \texttt{:id} & Retrieve current assigned mentor \\
\texttt{GET} & /api/students/:id/meetings & Student & URL param \texttt{:id} & Retrieve student meeting history \\
\midrule
\texttt{POST} & /api/meetings & Student & \texttt{\{mentor, student, date\}} & Book new conflict-free meeting \\
\texttt{GET} & /api/meetings & Protected & Filter query params & Query meetings by role/date \\
\texttt{PUT} & /api/meetings/:id/notes & Faculty & \texttt{\{notes\}} & Append faculty meeting conclusion \\
\texttt{PUT} & /api/meetings/:id/cancel & Protected & None & Cancel scheduled meeting \\
\midrule
\texttt{POST} & /api/notes & Faculty & \texttt{\{meetingId, content\}} & Create formal log note \\
\texttt{GET} & /api/notes/meeting/:id & Protected & URL param \texttt{:id} & Fetch notes for a meeting \\
\midrule
\texttt{POST} & /api/assignments & Admin & \texttt{\{studentId, facultyId\}} & Allocate mentor to student \\
\texttt{GET} & /api/assignments & Admin & None & Query all active assignments \\
\texttt{PUT} & /api/assignments/:id/reassign & Admin & \texttt{\{facultyId\}} & Reassign mentee to new faculty \\
\midrule
\texttt{POST} & /api/feedback & Student & \texttt{\{meeting, rating, comm\}} & Submit meeting satisfaction rating \\
\texttt{GET} & /api/feedback & Admin & None & Retrieve feedback aggregates \\
\bottomrule
\end{tabularx}
\caption{Complete RESTful API Endpoint Directory}
\label{tab:api}
\end{table}

\subsection{Standardized Response Structures}
All responses conform to a standardized JSON schema, ensuring consistent error parsing on the client:
\begin{lstlisting}[language=bash, caption={Sample Successful JSON Response Structure}]
{
  "success": true,
  "data": {
    "_id": "66e5a31bc9d4...",
    "student": "66e5a20...",
    "mentor": "66e5a15...",
    "date": "2026-09-20T10:30:00.000Z",
    "status": "scheduled",
    "agenda": "Review semester project design diagrams"
  }
}
\end{lstlisting}

% ============================================================================
% SECTION 7: FRONTEND COMPONENT ARCHITECTURE & USER EXPERIENCE
% ============================================================================
\section{Frontend Component Architecture \& User Experience}

\subsection{Component Topology and Route Protection}
The frontend is constructed around a declarative routing topology governed by React Router v6. Figure~\ref{fig:fe_tree} illustrates the component hierarchy.

\begin{figure}[H]
\centering
\small
\begin{tabular}{|p{0.96\textwidth}|}
\hline
\rowcolor{primary!10} \centering\textbf{\color{primary}Root Application Architecture: \texttt{App.jsx} $\&$ \texttt{AuthContext.jsx}} \tabularnewline
\hline
\vspace{0.1cm}
\begin{tabular}{@{}p{0.32\textwidth}p{0.32\textwidth}p{0.32\textwidth}@{}}
\textbf{\color{secondary}1. Public Auth Boundary} & \textbf{\color{secondary}2. Protected Route Gateways} & \textbf{\color{secondary}3. Shared UI Components} \\
$\bullet$ \texttt{/login} (\texttt{Login.jsx}) & $\bullet$ \texttt{ProtectedRoute [student]} & $\bullet$ \texttt{Navbar.jsx} (User State) \\
$\bullet$ \texttt{/register} (\texttt{Register.jsx}) & $\bullet$ \texttt{ProtectedRoute [faculty]} & $\bullet$ \texttt{Sidebar.jsx} (Role Links) \\
$\bullet$ Root redirect to \texttt{/login} & $\bullet$ \texttt{ProtectedRoute [admin]} & $\bullet$ \texttt{StatCard}, \texttt{Modal}, \texttt{EmptyState} \\
\end{tabular}
\vspace{0.1cm} \tabularnewline
\hline
\end{tabular}
\caption{Frontend Component Hierarchy \& Route Guard Architecture}
\label{fig:fe_tree}
\end{figure}

The \texttt{ProtectedRoute} component intercepts navigation requests, inspecting the active user's role claim within the React Context. Unauthorized route access triggers an immediate redirect to \texttt{/login} with state preservation.

\subsection{Walkthrough of Role-Based Interface Modules}

\subsubsection{1. Authentication Module}
Provides single-screen sign-in and account registration. The interface authenticates against \texttt{/api/users/login}, stores the returned JWT in \texttt{localStorage}, updates \texttt{AuthContext}, and automatically dispatches the user to their designated dashboard (\texttt{/student}, \texttt{/faculty}, or \texttt{/admin}).

\subsubsection{2. Student Portal Interface}
\begin{itemize}[leftmargin=1.4em, itemsep=0.2em]
  \item \textbf{Student Dashboard (\texttt{StudentDashboard.jsx}):} Features high-impact stat cards detailing assigned mentor name, total completed sessions, and upcoming appointments.
  \item \textbf{Assigned Mentor View (\texttt{MyMentor.jsx}):} Displays faculty contact information, department, office location, and areas of specialization.
  \item \textbf{Office-Hour Booking System (\texttt{BookMeeting.jsx}):} Displays open availability slots published by the mentor, enabling the student to select a time, state an agenda, and submit a booking request.
  \item \textbf{Meeting History \& Notes (\texttt{MeetingHistory.jsx}):} An auditable chronological table of past appointments, presenting faculty notes and allowing the student to submit a 1--5 star feedback score.
\end{itemize}

\subsubsection{3. Faculty Portal Interface}
\begin{itemize}[leftmargin=1.4em, itemsep=0.2em]
  \item \textbf{Faculty Dashboard (\texttt{FacultyDashboard.jsx}):} Overview of active advisee count, pending meeting requests, and today's schedule.
  \item \textbf{Mentees Directory (\texttt{Mentees.jsx}):} Tabular view of all assigned students, their roll numbers, academic focus areas, and cumulative meeting counts.
  \item \textbf{Availability Scheduler (\texttt{Availability.jsx}):} Intuitive controls for faculty to publish recurring or date-specific office-hour availability.
  \item \textbf{Meeting Manager \& Notes (\texttt{FacultyMeetings.jsx}, \texttt{AddNote.jsx}):} Enables faculty to mark sessions as completed or cancelled, and launch a note-taking modal to record confidential discussion points.
\end{itemize}

\subsubsection{4. Administrator Portal Interface}
\begin{itemize}[leftmargin=1.4em, itemsep=0.2em]
  \item \textbf{Admin Dashboard (\texttt{AdminDashboard.jsx}):} Institutional telemetry tracking aggregate student enrollment, mentor coverage percentage, and department-wide session volume.
  \item \textbf{Departmental Summary (\texttt{DepartmentSummary.jsx}):} Detailed breakdown of faculty advising loads, highlighting faculty members who exceed or fall below target mentee ratios.
  \item \textbf{Mentee Reassignment Engine (\texttt{ReassignMentee.jsx}):} Administrative intervention utility allowing rapid transfer of a student from one mentor to another.
\end{itemize}

% ============================================================================
% SECTION 8: KEY ENGINEERING IMPLEMENTATIONS
% ============================================================================
\section{Key Engineering Implementations}

\subsection{Conflict-Free Scheduling Algorithm}
To prevent schedule overlap and race conditions during simultaneous booking attempts, the backend validates slot availability prior to persisting appointment records. Formally, a booking conflict is triggered if there exists any meeting $M \in \text{Meetings}$ satisfying:
\begin{equation}
\begin{cases}
M.\text{mentor} = \text{TargetMentor} \\
M.\text{status} = \text{``scheduled''} \\
|M.\text{date} - \text{RequestedDate}| < \text{SlotDuration}
\end{cases}
\end{equation}
When a student initiates a booking request, the controller verifies both that the mentor is free during that time window and that the student does not have another active booking at that exact time.

\subsection{Asynchronous Error Handling Pipeline}
The application wraps all controller actions in an asynchronous error propagation utility (\texttt{asyncHandler.js}), obviating repetitive \texttt{try-catch} boilerplate and ensuring unhandled promise rejections are channeled into the centralized Express error middleware (\texttt{errorMiddleware.js}).

% ============================================================================
% SECTION 9: VERIFICATION, TESTING & QUALITY ASSURANCE
% ============================================================================
\section{Verification, Testing \& Quality Assurance}

\subsection{Acceptance Criteria Verification}
In accordance with the software engineering project specification, the prototype was evaluated against four core functional criteria:
\begin{enumerate}[leftmargin=1.4em, itemsep=0.2em]
  \item \textbf{Assignment Correctness (100\% Intra-Department):} Validated that 100\% of allocated students were matched exclusively with faculty belonging to their respective academic departments.
  \item \textbf{Zero Double-Booking (Concurrency Integrity):} Verified that the scheduling engine successfully rejects duplicate booking requests for the same faculty time slot.
  \item \textbf{Role Isolation \& Access Control:} Confirmed that student accounts attempting to invoke faculty or admin endpoints (e.g., \texttt{/api/assignments}) receive immediate HTTP 403 Forbidden rejections.
  \item \textbf{Administrative Visibility:} Ensured that department administrators can view live mentee distributions across faculty advisors without requiring manual reporting.
\end{enumerate}

\subsection{Prototype Test Suite Matrix}
Table~\ref{tab:tests} presents the execution results of the formal test suite executed on the prototype codebase.

\begin{table}[H]
\centering
\scriptsize
\begin{tabularx}{\textwidth}{@{}llp{3.5cm}Xc@{}}
\toprule
\textbf{Test ID} & \textbf{Target Module} & \textbf{Input Conditions} & \textbf{Expected Behavior} & \textbf{Status} \\
\midrule
\texttt{TC-01} & User Auth & Valid email \& password & HTTP 200, JWT token returned with correct role & \textbf{\color{successgreen}PASS} \\
\texttt{TC-02} & User Auth & Duplicate email registration & HTTP 400, ``User already exists'' error & \textbf{\color{successgreen}PASS} \\
\texttt{TC-03} & User Auth & Invalid password input & HTTP 401, ``Invalid credentials'' returned & \textbf{\color{successgreen}PASS} \\
\texttt{TC-04} & RBAC Guard & Student token accessing admin API & HTTP 403 Forbidden, request aborted & \textbf{\color{successgreen}PASS} \\
\texttt{TC-05} & Allocation & Assign student to mentor in CSE & Record linked; mentee count incremented & \textbf{\color{successgreen}PASS} \\
\texttt{TC-06} & Allocation & Reassign mentee to new mentor & Previous link marked inactive; new link active & \textbf{\color{successgreen}PASS} \\
\texttt{TC-07} & Availability & Faculty sets slots `Mon 10--12' & Availability string updated in Mentor model & \textbf{\color{successgreen}PASS} \\
\texttt{TC-08} & Booking & Student books open slot & HTTP 201, meeting created in `scheduled' state & \textbf{\color{successgreen}PASS} \\
\texttt{TC-09} & Concurrency & Concurrent booking of same slot & First succeeds (201); second rejected (409) & \textbf{\color{successgreen}PASS} \\
\texttt{TC-10} & Meeting Notes & Faculty appends meeting notes & Meeting updated; entry logged in Note DB & \textbf{\color{successgreen}PASS} \\
\texttt{TC-11} & Cancellation & Student cancels booked session & Meeting status transitions to `cancelled' & \textbf{\color{successgreen}PASS} \\
\texttt{TC-12} & Feedback & Student submits 5-star rating & Feedback record linked to meeting \& stored & \textbf{\color{successgreen}PASS} \\
\texttt{TC-13} & Dashboard & Admin fetches dept metrics & Aggregate counts match database records & \textbf{\color{successgreen}PASS} \\
\texttt{TC-14} & Health Check & \texttt{GET /api/health} & HTTP 200, system uptime confirmed & \textbf{\color{successgreen}PASS} \\
\bottomrule
\end{tabularx}
\caption{Prototype Stage Functional Verification Test Results}
\label{tab:tests}
\end{table}

% ============================================================================
% SECTION 10: PROJECT SCHEDULE & GANTT CHART TRACKING
% ============================================================================
\section{Project Schedule \& Gantt Chart Tracking}

\subsection{12-Week Semester Milestone Gantt Chart}
Figure~\ref{fig:gantt} illustrates the 12-week schedule for the UCS503P semester project, indicating completed phases, the active prototype review milestone, and upcoming pilot testing.

\begin{figure}[H]
\centering
\resizebox{\textwidth}{!}{%
\begin{ganttchart}[
    vgrid,
    hgrid,
    x unit=0.95cm,
    y unit chart=0.65cm,
    time slot format=simple,
    bar/.append style={fill=secondary!70, draw=primary},
    bar complete/.append style={fill=primary},
    group/.append style={fill=primary!30, draw=primary},
    milestone/.append style={fill=accent, draw=accent}
]{1}{12}
\gantttitle{UCS503P CMMS --- 12-Week Engineering Schedule}{12} \\
\gantttitlelist{1,...,12}{1} \\

\ganttgroup{Phase 1: Planning \& Setup}{1}{2} \\
\ganttbar[progress=100]{Requirements \& Architecture}{1}{2} \\
\ganttbar[progress=100]{Schema Design \& Repo Setup}{1}{2} \\

\ganttgroup{Phase 2: Core Modules}{3}{5} \\
\ganttbar[progress=100]{JWT Auth \& RBAC Engine}{3}{4} \\
\ganttbar[progress=100]{Mentee Allocation Engine}{4}{5} \\
\ganttbar[progress=100]{Base Frontend UI Setup}{3}{5} \\

\ganttgroup{Phase 3: Scheduling \& Notes}{6}{8} \\
\ganttbar[progress=100]{Office-Hours Booking API}{6}{7} \\
\ganttbar[progress=100]{Meeting Notes \& Logging}{7}{8} \\
\ganttbar[progress=100]{Scheduling UI Views}{6}{8} \\

\ganttgroup{Phase 4: Dashboards (Prototype)}{9}{10} \\
\ganttbar[progress=100]{Student \& Faculty Dashboards}{9}{10} \\
\ganttbar[progress=100]{Admin Analytics Dashboard}{9}{10} \\
\ganttmilestone{Prototype Evaluation Stage}{10} \\

\ganttgroup{Phase 5: Pilot Testing \& Polish}{11}{12} \\
\ganttbar[progress=25]{Department Pilot Run}{11}{11} \\
\ganttbar[progress=0]{Bug Fixes \& Final Defense}{12}{12}
\end{ganttchart}%
}
\caption{Project Gantt Chart Tracking Semester Progress (Prototype Delivered at W10)}
\label{fig:gantt}
\end{figure}

\subsection{Team Responsibilities \& Work Breakdown Matrix}
Table~\ref{tab:wbs} details the specific contributions of each team member towards delivering the prototype.

\begin{table}[H]
\centering
\small
\begin{tabularx}{\textwidth}{lp{3.5cm}X}
\toprule
\textbf{Member Name} & \textbf{Assigned Domain} & \textbf{Key Prototype Deliverables} \\
\midrule
\textbf{Aaditya Bansal} \newline (1024030768) & Backend Architecture \& System Integration Lead & Engineered Express 5 REST API gateway, JWT authentication pipeline, bcrypt password hashing, mentee allocation algorithm, and conflict-free booking endpoints. Integrated full-stack communications. \\
\midrule
\textbf{Jessica} \newline (1024030761) & Frontend Architecture \& UI/UX Design Lead & Developed modular React 18 SPA with Vite, designed the cohesive design system, implemented role-protected routing (\texttt{ProtectedRoute}), and authored all 12 responsive views for Student, Faculty, and Admin portals. \\
\midrule
\textbf{Harshveer Singh} \newline (1024030776) & Database Architecture \& Conflict Validation Lead & Architected MongoDB Atlas document schemas and Mongoose models (\texttt{User}, \texttt{Student}, \texttt{Mentor}, \texttt{Meeting}, \texttt{Note}, \texttt{Assignment}, \texttt{Feedback}). Formulated concurrency conflict resolution queries and seeded pilot data. \\
\bottomrule
\end{tabularx}
\caption{Work Breakdown and Individual Responsibility Matrix}
\label{tab:wbs}
\end{table}

% ============================================================================
% SECTION 11: EVALUATION AGAINST COURSE HEURISTICS
% ============================================================================
\section{Evaluation Against Course Heuristics}
The prototype directly fulfills the evaluation heuristics articulated in the UCS503P project criteria:
\begin{itemize}[leftmargin=1.4em, itemsep=0.3em]
  \item \textbf{1. Contextual Relevance:} The problem of mentorship opacity directly impacts students and faculty in the Computer Science and Engineering department at TIET, providing immediate institutional applicability.
  \item \textbf{2. Engine Availability:} Rather than developing proprietary database engines or novel cryptographic protocols, the system leverages battle-tested, mature industry engines (React, Node.js, Express, MongoDB, Bcrypt, JWT), allowing the engineering effort to focus on software architecture and usability.
  \item \textbf{3. Scalability:} Stateless JWT authentication and NoSQL document sharding capabilities in MongoDB Atlas ensure the platform can seamlessly scale from a single department prototype to campus-wide deployment across multiple faculties.
  \item \textbf{4. Measurable Evaluation Criteria:} All system criteria are quantifiable: 100\% allocation integrity, zero scheduling conflicts observed during concurrency tests, and sub-200ms round-trip API latency.
  \item \textbf{5. Time-to-Value \& Rapid Iteration:} In alignment with the two-week iteration heuristic, working increments were delivered across consecutive milestones, culminating in a fully demonstrable prototype.
\end{itemize}

% ============================================================================
% SECTION 12: LIMITATIONS, LESSONS LEARNED & ROADMAP
% ============================================================================
\section{Limitations, Lessons Learned \& Roadmap}

\subsection{Current Prototype Limitations}
While the prototype demonstrates complete functional coverage of the core mentorship workflow, several enhancements are scheduled for the final deployment phase:
\begin{itemize}[leftmargin=1.4em, itemsep=0.2em]
  \item \textbf{Notification Infrastructure:} The current prototype uses synchronous in-app UI notifications. Direct SMTP email and SMS alerts upon meeting confirmation will be integrated during Phase 5.
  \item \textbf{External Calendar Synchronization:} Automated two-way synchronization with Google Calendar and Microsoft Outlook via iCal webhooks is currently omitted to prevent external API rate-limiting during prototype grading.
  \item \textbf{Dynamic Load Balancing Heuristic:} While department admins can manually reassign mentees to balance faculty workloads, an automated constraint-satisfaction algorithm will be evaluated for the final release.
\end{itemize}

\subsection{Engineering Lessons Learned}
Developing the prototype revealed crucial full-stack engineering insights:
\begin{itemize}[leftmargin=1.4em, itemsep=0.2em]
  \item \textbf{Schema Normalization vs. Embedding:} Careful trade-offs were required between embedding meeting notes directly inside the \texttt{Meeting} document versus maintaining a discrete \texttt{Note} collection. Using a hybrid approach preserved fast dashboard queries while maintaining detailed historical notes.
  \item \textbf{Role State Hydration in React:} Ensuring that client-side page reloads did not cause momentary flashes or premature auth redirects required resilient token hydration within the \texttt{AuthContext} state provider.
\end{itemize}

\subsection{Roadmap Toward Final Submission}
\begin{enumerate}[leftmargin=1.4em, itemsep=0.2em]
  \item \textbf{Week 11 (Pilot Department Run):} Deploy prototype to a pilot group of 20 CSE students and 4 faculty members to collect user satisfaction metrics and monitor server telemetry under live conditions.
  \item \textbf{Week 12 (Final Polish \& Defense):} Implement UI responsiveness adjustments based on pilot user feedback, execute final code freeze, and prepare the comprehensive Final Project Report and live demonstration.
\end{enumerate}

% ============================================================================
% SECTION 13: CONCLUSION
% ============================================================================
\section{Conclusion}
The Prototype Stage of the \textbf{College Mentorship Management System (CMMS)} has successfully delivered a fully operational, multi-tier web platform that eliminates the chaos and opacity of traditional academic advising. By uniting student self-service booking, faculty office-hour management, timestamped interaction logging, and departmental administrative oversight into a cohesive, secure MERN application, CMMS provides a scalable foundation for modern collegiate mentorship. Comprehensive testing confirms that all prototype acceptance criteria have been achieved, positioning the project on schedule for its departmental pilot study and final submission.

% ============================================================================
% REFERENCES
% ============================================================================
\begin{thebibliography}{99}
\bibitem{pressman}
R.~S.~Pressman and B.~R.~Maxim, \emph{Software Engineering: A Practitioner's Approach}, 9th~ed., McGraw-Hill Education, 2020.

\bibitem{sommerville}
I.~Sommerville, \emph{Software Engineering}, 10th~ed., Pearson, 2016.

\bibitem{mern}
E.~M.~Banker, P.~Bakkum, S.~Verch, D.~Garrett, and T.~Hawkins, \emph{MongoDB in Action: Covers MongoDB version 3.0}, Manning Publications, 2016.

\bibitem{react}
Facebook Open Source, ``React --- A JavaScript library for building user interfaces,'' 2026. [Online]. Available: \url{https://react.dev/}.

\bibitem{express}
OpenJS Foundation, ``Express.js --- Fast, unopinionated, minimalist web framework for Node.js,'' 2026. [Online]. Available: \url{https://expressjs.com/}.

\bibitem{jwt}
M.~Jones, J.~Bradley, and N.~Sakimura, ``JSON Web Token (JWT),'' RFC 7519, Internet Engineering Task Force (IETF), May 2015.

\bibitem{tiet_se}
R.~B.~Venkataramaiyer, ``UCS503: Software Engineering Project Guidelines and Heuristics,'' Thapar Institute of Engineering and Technology, 2026. [Online]. Available: \url{https://tiet-ucs503.github.io/ucs503p-202627odd-template/}.
\end{thebibliography}

\end{document}
